\chapter{Shares, Corporate Actions and Indices}\label{ch:m1:shares-corporate-actions-indices}

On Friday 7 June 2024 the shares of Nvidia closed at \$1\,208.88. On Monday
morning they opened near \$121. Nobody lost a cent: over the weekend each
holder had received nine additional shares for every one held. A chart of
the raw closing prices shows a 90\% collapse; a backtest run on that chart
discovers a wonderful short-selling strategy; an index that included the
stock at its new price without further thought would have fallen with it.
Shares are not static objects. They pay out, split, multiply, merge and
disappear, and every system that stores a price must know what happened to
the thing being priced. This chapter is about those events, the arithmetic
that undoes them, and the indices that must be maintained through all of
them.

\section{What a share is}

\begin{definition}[Share and market capitalisation]\label{def:m1:shares-corporate-actions-indices:share}
A \emph{share}\index{share} is a unit of ownership in a company, giving its
holder a proportional claim on what the company distributes and, usually, a
vote. With $N$ shares outstanding at price $P$, the
\emph{market capitalisation}\index{market capitalisation} is $NP$.
\end{definition}

\begin{definition}[Free float]\label{def:m1:shares-corporate-actions-indices:float}
The \emph{free float}\index{free float} of a company is the fraction of its
shares available to public investors, excluding holdings that are not for
sale in the ordinary course: founders and other strategic holders,
governments, cross-holdings, shares locked up after a listing.
\end{definition}

A company may have several \emph{classes} of shares with different votes or
different \omterm{def:m1:shares-corporate-actions-indices:dividend}{dividends}, each with its own price and ticker; the same share may
also be listed on several exchanges in several currencies
(\cref{ch:m1:delta-one-instruments}). ``The price of the company'' is
therefore a convention; the data structures of \cref{ch:m1:reading-market-data}
exist to pin it down.

\section{Dividends}

\begin{definition}[Dividend, record date and ex-dividend date]\label{def:m1:shares-corporate-actions-indices:dividend}
A \emph{dividend}\index{dividend} is a distribution of cash (or shares) by a
company to its shareholders, of a declared amount per share. It is paid to
whoever is on the register at the close of the
\emph{record date}\index{record date}. The
\emph{ex-dividend date}\index{ex-dividend date} is the first day on which a
buyer of the share is no longer entitled to the dividend: trades made on or
after it settle too late to reach the register in time.
\end{definition}

\begin{dated}{2026-09}{Ex-date and record date in the United States}\label{dat:m1:shares-corporate-actions-indices:exdate}
Since US settlement moved to one day on 28 May 2024, the \omterm{def:m1:shares-corporate-actions-indices:dividend}{ex-dividend date} of
an ordinary distribution is the \omterm{def:m1:shares-corporate-actions-indices:dividend}{record date} itself: exchange rules now say
so. Under two-day settlement the ex-date was one business day \emph{before}
the \omterm{def:m1:shares-corporate-actions-indices:dividend}{record date}.
\end{dated}

\begin{omfigure}
\begin{tikzpicture}[font=\small]
\draw[-{Stealth[length=2.2mm]},thick,omExo] (-0.4,0) -- (15,0);
\foreach \x/\top/\bot in {0.8/Declaration/{amount and dates announced},5.2/{Last cum day}/{buyers still get the dividend},8.6/{Ex-date}/{price opens lower by about $D$},12.6/{Payment}/{cash reaches holders}}{
  \draw[omExo,thick] (\x,-0.12) -- (\x,0.12);
  \node[font=\small\bfseries,anchor=south] at (\x,0.2) {\top};
  \node[font=\footnotesize,omExo,anchor=north,align=center,text width=3.1cm] at (\x,-0.22) {\bot};
}
\draw[omThm,thick,decorate,decoration={brace,amplitude=4pt}] (8.6,1.0) -- (12.6,1.0)
  node[midway,above=5pt,font=\footnotesize,omThm] {holder is owed $D$: a receivable, not a price};
\end{tikzpicture}

\omcaption{The dates of a \omterm{def:m1:shares-corporate-actions-indices:dividend}{dividend}. With one-day settlement the ex-date
coincides with the \omterm{def:m1:shares-corporate-actions-indices:dividend}{record date}. Between the ex-date and the payment the
holder's wealth is the share \emph{plus} the \omterm{def:m1:shares-corporate-actions-indices:dividend}{dividend} owed: a P\&L system
that forgets the second part shows a loss that did not happen
(\cref{met:m1:pnl-and-positions:explain}).}
\end{omfigure}

\begin{proposition}[The ex-date drop]\label{prop:m1:shares-corporate-actions-indices:drop}
Ignoring taxes and one night's interest, and with no news between the two
auctions, a share that closes \omterm{def:m1:shares-corporate-actions-indices:dividend}{cum-dividend} at $P$ opens \omterm{def:m1:shares-corporate-actions-indices:dividend}{ex-dividend} at $P -
D$.
\end{proposition}
\begin{proof}
A share bought at the cum close is worth, the next morning, one \omterm{def:m1:shares-corporate-actions-indices:dividend}{ex-dividend}
share plus a claim to $D$ payable in a few days. If it opened above $P - D$,
buying at the close and selling at the open would earn a riskless profit;
below, the reverse trade would.
\end{proof}

In practice the drop is a little less than $D$ where \omterm{def:m1:shares-corporate-actions-indices:dividend}{dividends} are taxed more
heavily than capital gains, because the marginal holder values a dollar of
\omterm{def:m1:shares-corporate-actions-indices:dividend}{dividend} at less than a dollar; the difference, and the lending of shares
across the ex-date that exploits it, return in
\cref{ch:m1:stock-loan-and-short-selling}.

\section{Splits, rights issues and spin-offs}

\begin{definition}[Stock split]\label{def:m1:shares-corporate-actions-indices:split}
In a $k$-for-1 \emph{stock split}\index{stock split} every share becomes $k$
shares. The number of shares is multiplied by $k$, the price is divided by
$k$, and nothing else changes. A \emph{reverse split} has $k<1$.
\end{definition}

\begin{definition}[Rights issue]\label{def:m1:shares-corporate-actions-indices:rights}
In a \emph{rights issue}\index{rights issue} a company raises capital by
giving its shareholders the right to buy $n$ new shares for every $N$ held,
at a subscription price $S$ below the market price. The rights are usually
tradable for a few weeks.
\end{definition}

\begin{proposition}[Theoretical ex-rights price]\label{prop:m1:shares-corporate-actions-indices:terp}
If the share closes at $P$ before an $n$-for-$N$ \omterm{def:m1:shares-corporate-actions-indices:rights}{rights issue} at $S$, its
theoretical price once the rights have detached is
\[
\mathrm{TERP} = \frac{NP + nS}{N+n},
\]
and one right to subscribe one new share is worth $\mathrm{TERP} - S$.
\end{proposition}
\begin{proof}
After the issue, $N+n$ shares own the old company, worth $NP$, plus $nS$ of
new cash. A right lets its owner pay $S$ for something worth
$\mathrm{TERP}$.
\end{proof}

\begin{example}[One for four at six]\label{ex:m1:shares-corporate-actions-indices:terp}
A share trades at 10.00; the company issues 1 new share for 4 held at 6.00.
$\mathrm{TERP} = (40 + 6)/5 = 9.20$. A right to one new share is worth
$3.20$; a holder of 4 shares receives one right and owns $4 \times 9.20 +
3.20 = 40.00$, as before. The share price falls 8\% on the ex-date and
nobody is poorer; a holder who ignores the rights, however, gives away 3.20.
\end{example}

\begin{definition}[Spin-off]\label{def:m1:shares-corporate-actions-indices:spinoff}
In a \emph{spin-off}\index{spin-off} a company distributes to its
shareholders the shares of a subsidiary, which becomes a separately listed
company. The parent's price falls on the ex-date by about the value
distributed per share.
\end{definition}

Mergers complete the list: in a cash merger the share disappears against a
payment; in a share merger it becomes a fixed number of the acquirer's
shares. For data purposes both end a price series, and the rule is the same
as everywhere in this chapter: the \emph{holder's wealth} is continuous; the
\emph{price} is not.

\section{Adjustment factors}

\begin{definition}[Adjustment factor]\label{def:m1:shares-corporate-actions-indices:factor}
The \emph{adjustment factor}\index{adjustment factor} of a corporate action
is the number $f$ by which every price before its ex-date is multiplied so
that the adjusted series has no jump caused by the action. A
\emph{back-adjusted} price is the raw price times the product of the factors
of all later actions.
\end{definition}

\begin{method}[The factor of each action]\label{met:m1:shares-corporate-actions-indices:factors}
With $P$ the last cum price:
\begin{enumerate}
\item $k$-for-1 split: $f = 1/k$ (volumes are multiplied by $k$).
\item Cash \omterm{def:m1:shares-corporate-actions-indices:dividend}{dividend} $D$: $f = (P-D)/P$.
\item \omterm{def:m1:shares-corporate-actions-indices:rights}{Rights issue}: $f = \mathrm{TERP}/P$.
\item \omterm{def:m1:shares-corporate-actions-indices:spinoff}{Spin-off} worth $V$ per parent share: $f = (P-V)/P$.
\end{enumerate}
Store the raw prices and the actions; compute adjusted series on demand.
Adjusted prices change every time a new action occurs, so a stored adjusted
price is a number with an unstated date.
\end{method}

\begin{proposition}[What the dividend factor does, exactly]\label{prop:m1:shares-corporate-actions-indices:exact}
Let the share close cum at $P$ and ex at $P_{\mathrm{ex}}$. The return of the
series adjusted with $f = (P-D)/P$ is $P_{\mathrm{ex}}/(P-D) - 1$, while the
holder's total return is $(P_{\mathrm{ex}}+D)/P - 1$. They are equal when
$P_{\mathrm{ex}} = P - D$ and differ at second order otherwise. The factor
$f^\star = P_{\mathrm{ex}}/(P_{\mathrm{ex}}+D)$, which assumes the \omterm{def:m1:shares-corporate-actions-indices:dividend}{dividend} is
reinvested at the ex price, reproduces the total return in every case.
\end{proposition}
\begin{proof}
The adjusted cum price is $fP = P - D$, giving the first return. With
$f^\star$ the adjusted cum price is $P P_{\mathrm{ex}}/(P_{\mathrm{ex}}+D)$ and
the return is $(P_{\mathrm{ex}}+D)/P - 1$. For the difference, write
$P_{\mathrm{ex}} = (P-D)(1+\varepsilon)$: the two returns are $\varepsilon$
and $\varepsilon\,(P-D)/P$, which differ by $\varepsilon D/P$.
\end{proof}

\begin{example}[Two basis points]\label{ex:m1:shares-corporate-actions-indices:approx}
$P = 102$, $D = 2$, and the stock closes ex at 99, a 1\% fall net of the
\omterm{def:m1:shares-corporate-actions-indices:dividend}{dividend}. The standard factor gives $-1.000\%$; the holder earned $101/102 -
1 = -0.980\%$. The error, two basis points, is $\varepsilon D/P$. Vendors
differ in which factor they publish: two ``adjusted close'' series for the
same stock need not agree, and a research database must state its
convention.
\end{example}

\begin{omfigure}
\begin{tikzpicture}
\begin{axis}[width=11cm,height=5.2cm,ymode=log,
  xlabel={trading day},ylabel={price (log scale)},
  tick label style={font=\small},label style={font=\small},
  xmin=0,xmax=250,ymin=25,ymax=900,grid=major,grid style={gray!25},
  ytick={30,50,100,200,400,800},yticklabels={30,50,100,200,400,800},
  legend columns=2,legend style={font=\footnotesize,at={(0.5,-0.3)},anchor=north,draw=none,fill=none,/tikz/every even column/.append style={column sep=8pt}},
  table/col sep=comma]
\addplot[omThm,thick] table[x=day,y=raw]{figdata/markets-1/08-shares-corporate-actions-indices/split.csv};
\addplot[omDef,very thick] table[x=day,y=adjusted]{figdata/markets-1/08-shares-corporate-actions-indices/split.csv};
\legend{raw close,back-adjusted close}
\end{axis}
\end{tikzpicture}

\omcaption{A simulated stock with a ten-for-one split on day 150 and
\omterm{def:m1:shares-corporate-actions-indices:dividend}{dividends} on days 60 and 200. The raw series shows a 90\% fall that no
holder experienced; the back-adjusted series is continuous, and \emph{all}
its values before day 200 differ from the prices that actually traded. Data:
the chapter's script.}\label{fig:m1:shares-corporate-actions-indices:split}
\end{omfigure}

\begin{remark}[Three uses, three series]\label{rem:m1:shares-corporate-actions-indices:uses}
Use \emph{raw} prices for anything that touches an order: limit prices, tick
sizes, fills, fees, the simulator. Use \emph{adjusted} prices for returns,
volatilities and signals computed across an ex-date. And in a backtest, use
the adjusted series \emph{as it was known on each date}: a signal computed
today from back-adjusted prices of 2019 has seen the split of 2024. That
discipline, called point-in-time data, is the subject of One Quant Book 7.
\end{remark}

\section{Indices}

An index is a portfolio defined by a rule, whose value is published as a
single number. Two design choices define it: how constituents are weighted,
and how the number is kept continuous when the portfolio changes.

\begin{definition}[Index divisor]\label{def:m1:shares-corporate-actions-indices:divisor}
A capitalisation-weighted index with constituents $i$ of price $p_i$, shares
outstanding $s_i$ and float factor $\phi_i \in [0,1]$ has level
\[
I = \frac{1}{\Delta}\sum_i p_i\,s_i\,\phi_i ,
\]
where the \emph{index divisor}\index{index divisor} $\Delta$ is chosen at
inception to give a round starting level and is thereafter changed only to
offset changes in the numerator that are not price moves.
\end{definition}

\begin{proposition}[Divisor adjustment]\label{prop:m1:shares-corporate-actions-indices:divisoradj}
When, at given prices, a change of constituents, shares or float moves the
index market value from $M$ to $M'$, the level is unchanged if and only if
the divisor becomes $\Delta' = \Delta\,M'/M$.
\end{proposition}
\begin{proof}
$I = M/\Delta = M'/\Delta'$.
\end{proof}

A split changes $p_i$ and $s_i$ in opposite proportions and leaves $M$
unchanged: a capitalisation-weighted index needs no adjustment for it. A
\emph{price-weighted} index, $I = \sum_i p_i/\Delta$, does: the split
constituent's weight falls tenfold and the divisor must shrink to
compensate. The best-known price-weighted index, the Dow Jones Industrial
Average, has been adjusted so many times that its divisor is well below
one: a one-dollar move in any of its thirty stocks moves the index by
several points.

\begin{definition}[Total return index]\label{def:m1:shares-corporate-actions-indices:tri}
A \emph{price index} reflects only price changes: on an ex-date it falls
with its constituent. A \emph{total return index}\index{total return index}
reinvests every \omterm{def:m1:shares-corporate-actions-indices:dividend}{dividend} in the index on its ex-date (a \emph{net} total
return index after deducting a withholding tax), and is the correct benchmark
for a fund that receives the \omterm{def:m1:shares-corporate-actions-indices:dividend}{dividends}.
\end{definition}

\begin{omfigure}
\begin{tikzpicture}
\begin{axis}[width=11cm,height=4.8cm,ybar,bar width=7pt,
  ylabel={weight (\%)},symbolic x coords={A,B,C,D,E},xtick=data,
  tick label style={font=\small},label style={font=\small},
  ymin=0,ymax=70,ymajorgrids,grid style={gray!25},enlarge x limits=0.15,area legend,
  legend columns=3,legend style={font=\footnotesize,at={(0.5,-0.2)},anchor=north,draw=none,fill=none,/tikz/every even column/.append style={column sep=8pt}},
  table/col sep=comma]
\addplot[fill=omDef!70,draw=omDef] table[x=stock,y=cap]{figdata/markets-1/08-shares-corporate-actions-indices/weights.csv};
\addplot[fill=omThm!55,draw=omThm] table[x=stock,y=price]{figdata/markets-1/08-shares-corporate-actions-indices/weights.csv};
\addplot[fill=omMeth!55,draw=omMeth] table[x=stock,y=equal]{figdata/markets-1/08-shares-corporate-actions-indices/weights.csv};
\legend{float-adjusted capitalisation,price,equal}
\end{axis}
\end{tikzpicture}

\omcaption{Five illustrative stocks under three weighting rules. Stock D, the
smallest company, dominates the price-weighted index because its share price
is 900; stock B, one of the largest companies, nearly vanishes from it. Data: the chapter's script.}
\end{omfigure}

\begin{omfigure}
\begin{tikzpicture}
\begin{axis}[width=11cm,height=4.6cm,
  xlabel={trading day},ylabel={index level},
  tick label style={font=\small},label style={font=\small},
  xmin=0,xmax=59,ymin=900,ymax=1400,grid=major,grid style={gray!25},
  yticklabel style={/pgf/number format/1000 sep={\,}},
  legend columns=2,legend style={font=\footnotesize,at={(0.5,-0.34)},anchor=north,draw=none,fill=none,/tikz/every even column/.append style={column sep=8pt}},
  table/col sep=comma]
\addplot[omDef,very thick] table[x=day,y=maintained]{figdata/markets-1/08-shares-corporate-actions-indices/rebase.csv};
\addplot[omThm,thick,dashed] table[x=day,y=naive]{figdata/markets-1/08-shares-corporate-actions-indices/rebase.csv};
\legend{divisor adjusted on day 30,divisor left unchanged}
\end{axis}
\end{tikzpicture}

\omcaption{On day 30 stock B leaves a five-stock index and a larger company
enters. With \cref{prop:m1:shares-corporate-actions-indices:divisoradj} the
level is continuous; without it the index ``gains'' a third in a day in which
no price moved. Data: the chapter's script.}
\end{omfigure}

\section{Tutorial: adjusting a series and maintaining an index}\label{tut:m1:shares-corporate-actions-indices:adjust}

\begin{tutorial}
\textbf{Goal.} Back-adjust a price series through a split and two \omterm{def:m1:shares-corporate-actions-indices:dividend}{dividends},
then carry an index through a constituent change. \textbf{End state:}
\cref{fig:m1:shares-corporate-actions-indices:split} and the two index paths
of the last figure.
\begin{tutsteps}
\item \textbf{Factors and back-adjustment.} Each action multiplies every
\emph{earlier} price by its factor.
\omcode{code/markets-1/08-shares-corporate-actions-indices/python/corpact.py}{7}{14}{Split and dividend factors.}\label{lst:m1:shares-corporate-actions-indices:factors}
\omcode{code/markets-1/08-shares-corporate-actions-indices/python/corpact.py}{31}{37}{Back-adjustment: one line per action.}\label{lst:m1:shares-corporate-actions-indices:backadj}
\item \textbf{Check.} The figure script builds the raw series from a known
total-return path and asserts that the adjusted series has exactly that
path's returns.
\item \textbf{The index.} The divisor moves only inside \texttt{rebase}, and
only by the ratio of market values at unchanged prices.
\omcode{code/markets-1/08-shares-corporate-actions-indices/python/corpact.py}{41}{58}{A float-adjusted capitalisation-weighted index with divisor maintenance.}\label{lst:m1:shares-corporate-actions-indices:index}
\item \textbf{Break it.} Replace the constituent without calling
\texttt{rebase}: the level jumps from 981 to 1\,320.
\end{tutsteps}
\textbf{What to change next.} Add a total-return version of the index: on
each ex-date increase the level by the \omterm{def:m1:shares-corporate-actions-indices:dividend}{dividend}'s index points. Then apply
the ten-for-one split to a price-weighted version and compute the new
divisor by hand.
\end{tutorial}

\section{Build: the corporate-action adjuster}\label{bld:m1:shares-corporate-actions-indices:adjuster}

\begin{build}
\bfield{Purpose} One service in the miniature firm answers ``what is the
adjusted price of this symbol on this date, as known on that date?''.
Research, risk and the P\&L keeper all call it; none of them stores adjusted
prices.
\bfield{Interface} \texttt{Action(symbol, ex\_date, kind, params,
announced)} with \texttt{kind} in \texttt{split}, \texttt{cash\_dividend},
\texttt{rights}, \texttt{spinoff}; \texttt{Adjuster.add(action)};
\texttt{factor(symbol, date, as\_of)} returning the product of the factors of
actions with \texttt{date < ex\_date <= as\_of}; \texttt{adjust\_position(
position, action)} for the keeper of \cref{ch:m1:pnl-and-positions}.
\bfield{Rules} Factors follow
\cref{met:m1:shares-corporate-actions-indices:factors}; the \omterm{def:m1:shares-corporate-actions-indices:dividend}{dividend}
convention is a constructor argument and is recorded in every output. An
action is invisible to queries whose \texttt{as\_of} precedes its
announcement. A split multiplies a position's quantity by $k$ and divides
its \omterm{def:m1:pnl-and-positions:realised}{average cost} by $k$, leaving cash untouched; a cash \omterm{def:m1:shares-corporate-actions-indices:dividend}{dividend} posts $qD$
to the ledger on the ex-date as a receivable.
\bfield{Acceptance tests} \texttt{code/firm/corpactions/tests/}: continuity
through a split; the as-of rule; position invariance (quantity times \omterm{def:m1:pnl-and-positions:realised}{average
cost} unchanged by a split); the one-for-four rights example.
\bfield{Stretch} Mergers: map a position in the target into cash or into
shares of the acquirer and close the price series.
\end{build}

\begin{omsources}
\item CNBC, ``Nvidia announces 10-for-1 stock split'', 22 May 2024; Yahoo
Finance, ``Nvidia stock rises after 10-for-1 stock split'', 10 June 2024.
\item US Securities and Exchange Commission, Release 34-99881 (NYSE Arca
rule change on ex-dividend dates under the one-day settlement cycle), 2024.
\item S\&P Dow Jones Indices, \emph{Index Mathematics Methodology}, April
2026; \emph{Dow Jones Averages Methodology}, May 2026; \emph{How the Dow
Works} (education paper).
\item E.~Elton and M.~Gruber, ``Marginal stockholder tax rates and the
clientele effect'', \emph{Review of Economics and Statistics} 52 (1970) ---
the classic study of the ex-date drop.
\end{omsources}

\section{Exercises}

\begin{exercise}[$\star$]\label{exo:m1:shares-corporate-actions-indices:1}
A company has 800 million shares at \$65 and a \omterm{def:m1:shares-corporate-actions-indices:float}{free float} of 70\%. Give its
\omterm{def:m1:shares-corporate-actions-indices:share}{market capitalisation} and its float-adjusted capitalisation.
\end{exercise}

\begin{exercise}[$\star$]\label{exo:m1:shares-corporate-actions-indices:2}
A stock closes at \$240 before a 3-for-2 split. Give the theoretical opening
price, the \omterm{def:m1:shares-corporate-actions-indices:factor}{adjustment factor}, and the new position of a holder of 500
shares.
\end{exercise}

\begin{exercise}[$\star$]\label{exo:m1:shares-corporate-actions-indices:3}
A stock closes cum at \$80.00 with a \omterm{def:m1:shares-corporate-actions-indices:dividend}{dividend} of \$1.20. Give the theoretical
ex price, the \omterm{def:m1:shares-corporate-actions-indices:factor}{adjustment factor}, and the back-adjusted value of a price of
\$76.00 observed a month earlier.
\end{exercise}

\begin{exercise}[$\star\star$]\label{exo:m1:shares-corporate-actions-indices:4}
A company at \euro 25 announces a 2-for-7 \omterm{def:m1:shares-corporate-actions-indices:rights}{rights issue} at \euro 16. Compute
the TERP, the value of one right, the \omterm{def:m1:shares-corporate-actions-indices:factor}{adjustment factor}, and check that a
holder of 700 shares who sells all the rights is as wealthy as before.
\end{exercise}

\begin{exercise}[$\star\star$]\label{exo:m1:shares-corporate-actions-indices:5}
A price-weighted index of three stocks at \$50, \$120 and \$330 has a divisor
of 0.5. Give its level. The third \omterm{def:m1:shares-corporate-actions-indices:split}{stock splits} 3-for-1: give the new divisor
and the new weight of each stock.
\end{exercise}

\begin{exercise}[$\star\star$]\label{exo:m1:shares-corporate-actions-indices:6}
A stock closes cum at \$50.00 with a \omterm{def:m1:shares-corporate-actions-indices:dividend}{dividend} of \$2.50 and closes on the
ex-date at \$48.45. Compute the ex-date return with the standard factor, the
holder's total return, and the difference. Check it against the formula of
\cref{prop:m1:shares-corporate-actions-indices:exact}.
\end{exercise}

\begin{exercise}[$\star\star\star$]\label{exo:m1:shares-corporate-actions-indices:7}
\emph{Coding.} A capitalisation-weighted index has two stocks: X with 100
million shares at \$20 and Y with 50 million at \$60, floats of 1, and an
initial level of 1\,000. Using \texttt{CapIndex}: (a) give the divisor; (b) Y
buys back 10 million shares at unchanged prices: give the new divisor; (c)
prices then move to \$22 and \$57: give the level.
\end{exercise}

\begin{exercise}[$\star\star\star$]\label{exo:m1:shares-corporate-actions-indices:8}
\emph{Find the flaw.} A researcher downloads today's back-adjusted daily
closes for 3\,000 stocks over twenty years and tests the rule ``buy stocks
whose price is below \$5''. The backtest is excellent. Identify two distinct
ways in which the adjusted prices have contaminated the test.
\end{exercise}

\section{Problem: Maintaining an Index Through a Bad Week}

\begin{problem}[{Weekend problem --- five days at an index provider}]\label{pb:m1:shares-corporate-actions-indices:1}
You maintain a float-adjusted capitalisation-weighted index of four stocks.
On Friday's close:
\begin{center}\small
\begin{tabular}{@{}lrrr@{}}
\hline
Stock & Price & Shares (millions) & Float \\
\hline
P & 40.00 & 500 & 1.00 \\
Q & 150.00 & 200 & 0.80 \\
R & 12.00 & 1\,000 & 0.50 \\
S & 75.00 & 400 & 0.90 \\
\hline
\end{tabular}
\end{center}
The index closed at 2\,500.00.

\medskip\noindent\textbf{Part I --- Friday.}
\begin{enumerate}
\item Compute each stock's float-adjusted capitalisation and the index
market value.
\item Give the divisor.
\item Give each stock's weight.
\item By how many index points does the index move if Q rises 1\%?
\end{enumerate}

\medskip\noindent\textbf{Part II --- Monday: a split and a \omterm{def:m1:shares-corporate-actions-indices:dividend}{dividend}.}
\begin{enumerate}[resume]
\item Q splits 3-for-1 before the open. What happens to its price, its
shares and the divisor?
\item P goes ex a \omterm{def:m1:shares-corporate-actions-indices:dividend}{dividend} of 1.00 and opens at its theoretical ex price;
nothing else moves. Give the price index. Is the divisor adjusted?
\item Give the level of the \emph{total return} version of the index, which
also closed Friday at 2\,500.00.
\item A fund tracking the price index received the \omterm{def:m1:shares-corporate-actions-indices:dividend}{dividend}. Which index
should it be compared with?
\end{enumerate}

\medskip\noindent\textbf{Part III --- Wednesday: a \omterm{def:m1:shares-corporate-actions-indices:rights}{rights issue}.}
Prices are those of Monday's open.
\begin{enumerate}[resume]
\item R announces a 1-for-4 \omterm{def:m1:shares-corporate-actions-indices:rights}{rights issue} at 8.00, effective Wednesday. Give
the TERP and R's new share count.
\item Treating the issue as fully subscribed, give R's new float-adjusted
capitalisation at the TERP and the new index market value.
\item Give the new divisor.
\item A tracking fund holds R at index weight. What must it do, in words,
to stay in line?
\end{enumerate}

\medskip\noindent\textbf{Part IV --- Friday: a replacement.}
Prices are unchanged since Wednesday.
\begin{enumerate}[resume]
\item S is acquired for cash and leaves the index at 75.00. It is replaced
by T: price 30.00, 600 million shares, float 0.75. Give the new market value
and divisor.
\item Give the new weights.
\item A fund with \$2 billion tracking the index must sell its S and buy T.
How many dollars of each?
\item T trades \$150 million a day. Express the purchase in days of volume
and estimate its cost with \cref{met:m1:the-sell-side:sqrt} for a daily
volatility of 2\%.
\item Who is on the other side of that purchase, and when did they buy?
\item The index provider announced the change five days earlier. Why not
announce and implement on the same day?
\item State the \emph{named result}: the \omterm{def:m1:shares-corporate-actions-indices:divisor}{index divisor} at the end of the
week.
\item In one sentence: what is the one thing a divisor adjustment must never
do?
\end{enumerate}
\end{problem}

\section{Interview questions}

\begin{interviewq}[$\star$ \iqroles{researcher, developer, trader}]\label{iq:m1:shares-corporate-actions-indices:1}
A stock's price drops 50\% overnight and the company's value is unchanged.
What happened, and what must your data pipeline do about it?
\end{interviewq}

\begin{interviewq}[$\star$ \iqroles{trader, researcher}]\label{iq:m1:shares-corporate-actions-indices:2}
By how much should a stock fall on its \omterm{def:m1:shares-corporate-actions-indices:dividend}{ex-dividend date}? Why might it fall by
less?
\end{interviewq}

\begin{interviewq}[$\star\star$ \iqroles{researcher, mle}]\label{iq:m1:shares-corporate-actions-indices:3}
When should you use raw prices and when adjusted prices? Give one bug caused
by each wrong choice.
\end{interviewq}

\begin{interviewq}[$\star\star$ \iqroles{trader, researcher}]\label{iq:m1:shares-corporate-actions-indices:4}
Why does a \omterm{def:m1:shares-corporate-actions-indices:split}{stock split} change the weights of a price-weighted index and not
those of a capitalisation-weighted one?
\end{interviewq}

\begin{interviewq}[$\star\star$ \iqroles{developer}]\label{iq:m1:shares-corporate-actions-indices:5}
Design the storage of prices and corporate actions for a research database.
What do you store, what do you compute on read, and why?
\end{interviewq}

\begin{interviewq}[$\star\star\star$ \iqroles{trader, researcher}]\label{iq:m1:shares-corporate-actions-indices:6}
A company announces a 1-for-3 \omterm{def:m1:shares-corporate-actions-indices:rights}{rights issue} at a 40\% discount. The shares
fall 12\% on the announcement. Is that ``because of the dilution''?
\end{interviewq}
